\newpage
\begin{center}
  \textbf{\large ЗАКЛЮЧЕНИЕ}
\end{center}
\addcontentsline{toc}{chapter}{ЗАКЛЮЧЕНИЕ}

В работе была рассмотрена задача реконструкции бинарных визуальных стимулов
размером $6 \times 6$ пикселей по многоканальным EEG-сигналам, записанным в
фазе воспоминания ранее предъявленного изображения. В отличие от классической
классификации стимулов, такая постановка требует восстановления локальной
пиксельной структуры образа и поэтому была сформулирована как многовыходная
бинарная задача с 36 целевыми переменными.

Для решения этой задачи был реализован воспроизводимый экспериментальный
пайплайн. Использовался готовый предварительно обработанный EEG/EOG-набор
данных, из которого для основных экспериментов была выделена подвыборка
случайных визуальных стимулов фазы воспоминания. Для классических моделей
были построены временные, спектральные, пространственные и локально-паттерновые
признаки. Для нейросетевых моделей были подготовлены спектральные входы
FFT/Welch, Morlet, Superlet и STFT. В сравнение вошли референсная Logistic
Regression модель, набор классических моделей машинного обучения и компактные
сверточные архитектуры EEGNet, DeepConvNet и ShallowConvNet.

Оценка проводилась в двух независимых по смыслу протоколах. Межсубъектный
протокол проверял перенос на новых испытуемых, а двунаправленный межпробный
протокол проверял перенос между экспериментальными пробами у испытуемых, для
которых были доступны обе пробы. Во всех случаях сохранялся контроль утечек
между обучением и тестом: разбиения учитывали испытуемых, пробы, ключи
примеров, random seed и полные целевые изображения. Для интерпретации
результатов использовались попиксельные метрики, субъектно-кластерный
bootstrap и парное bootstrap-сравнение с поправкой Holm.

Главный экспериментальный вывод является консервативным. В межсубъектном
протоколе описательным лидером среди обучаемых моделей стала независимая
ридж-регрессия со сбалансированной точностью 0.518382. В двунаправленном
межпробном протоколе описательным лидером стала модель DeepConvNet на
STFT-входе со средней сбалансированной точностью 0.512011. Однако обе оценки
остаются близкими к случайному уровню 0.5, а финальное парное
bootstrap-сравнение не выявило статистически надежного превосходства над
логистической регрессией. Минимальное скорректированное по Holm значение
$p$ в финальном сравнении составило 0.273000. Доля полных совпадений для
изображения $6 \times 6$ была равна нулю для всех обучаемых моделей
финального сравнения.

Таким образом, работа не подтверждает надежную декодируемость случайного
визуального образа из EEG в текущей постановке и на текущем объеме данных.
При этом был получен важный инженерный результат: построена проверяемая
экспериментальная инфраструктура, включающая воспроизводимые признаки и
спектральные представления, train-only нормализацию и подбор параметров,
сохранение неизменяемых артефактов, аудит утечек, bootstrap-оценку
неопределенности, безопасную проверку результатов и программно сгенерированные
иллюстрации для анализа.

Ограничения работы связаны прежде всего с малым числом экспериментальных
примеров, отсутствием внешней валидации на независимом корпусе и высокой
строгостью полной реконструкции 36 бинарных пикселей. Кроме того, EOG-записи
в данной работе использовались как часть исходного экспериментального
контекста и контроля качества, но не включались в модели как дополнительный
вход. Нейросетевые архитектуры обучались на малом объеме данных, поэтому их
результаты следует рассматривать как проверку реализуемости пайплайна, а не
как окончательное доказательство преимущества глубоких моделей.

Дальнейшее развитие работы может включать расширение числа испытуемых и проб,
независимую внешнюю валидацию, персональную адаптацию моделей, использование
EOG или других модальностей как дополнительных входов, анализ более богатых
стимульных наборов, а также проверку методов переноса обучения и
предобученных EEG-представлений. При этом ключевым требованием остается
строгое разделение обучающих и тестовых данных и сохранение субъектного
уровня анализа при утверждениях о переносимости модели.
